\documentclass[11pt]{article}
\usepackage[a4paper,margin=2.5cm]{geometry}
\usepackage{amsmath}
\usepackage{amssymb}
\usepackage{enumitem}
\usepackage{parskip}
\setlist[enumerate,2]{label=(\alph*), itemsep=1pt, topsep=2pt}
\setlength{\parindent}{0pt}
\newcommand{\correct}[1]{\textbf{#1}\;$\checkmark$}
\begin{document}
\section*{Repeated games}
\begin{enumerate}[itemsep=6pt]
  \item In an infinitely repeated game, the discount factor \(\delta\) can be read as ...
  \begin{enumerate}
    \item the probability that the game continues to another round
    \item the number of rounds already played
    \item the number of players
    \item the payoff from defecting
  \end{enumerate}
  \item The Grudger strategy ...
  \begin{enumerate}
    \item copies the opponent's last move
    \item always defects
    \item cooperates until the opponent defects, then defects for ever
    \item cooperates with probability \(1/2\)
  \end{enumerate}
  \item The Folk Theorem says that, for patient enough players, ...
  \begin{enumerate}
    \item cooperation is impossible
    \item any individually rational payoff pair can be sustained in a subgame perfect equilibrium
    \item the game must be zero sum
    \item only the stage Nash equilibrium can occur
  \end{enumerate}
  \item In the finitely repeated Prisoner's Dilemma with a known last round, ...
  \begin{enumerate}
    \item cooperation is sustained throughout
    \item the players cooperate only in the last round
    \item there is no subgame perfect equilibrium
    \item backward induction forces defection in every round
  \end{enumerate}
  \item The individually rational payoff of a player is ...
  \begin{enumerate}
    \item the largest payoff in the game
    \item the most they can guarantee themselves whatever the opponent does (their minimax value)
    \item always zero
    \item the payoff from mutual cooperation
  \end{enumerate}
  \item As the temptation payoff \(T\) increases, the discount factor needed to sustain cooperation ...
  \begin{enumerate}
    \item stays the same
    \item decreases
    \item becomes negative
    \item increases (players must be more patient)
  \end{enumerate}
  \item A strategy in a repeated game is ...
  \begin{enumerate}
    \item the payoff matrix
    \item the discount factor
    \item a single action played every round
    \item a mapping from the history of play to an action in the stage game
  \end{enumerate}
  \item The discounted value of receiving a payoff \(c\) every round is ...
  \begin{enumerate}
    \item \(c \delta\)
    \item \(c / \delta\)
    \item \(c / (1 - \delta)\)
    \item \(c (1 - \delta)\)
  \end{enumerate}
  \item The average utility of a payoff stream is defined as ...
  \begin{enumerate}
    \item \((1 - \delta)\) times the discounted total
    \item the largest stage payoff
    \item the discounted total
    \item the payoff in the first round
  \end{enumerate}
  \item Playing a Nash equilibrium of the stage game in every round is ...
  \begin{enumerate}
    \item always a subgame perfect equilibrium of the repeated game
    \item impossible to compute
    \item never an equilibrium
    \item only a Nash equilibrium if the game is finite
  \end{enumerate}
  \item Cooperation can be sustained in the infinitely repeated Prisoner's Dilemma because ...
  \begin{enumerate}
    \item the stage game has no Nash equilibrium
    \item players cannot remember the past
    \item defection is impossible
    \item there is no last round, so future punishment can deter defection
  \end{enumerate}
  \item The expected number of rounds when \(\delta\) is the continuation probability is ...
  \begin{enumerate}
    \item infinite for any \(\delta\)
    \item \(1 - \delta\)
    \item \(1 / (1 - \delta)\)
    \item \(\delta\)
  \end{enumerate}
  \item In Grudger, the punishment phase consists of ...
  \begin{enumerate}
    \item returning to cooperation after one round
    \item alternating cooperation and defection
    \item both players playing the stage Nash equilibrium (mutual defection) for ever
    \item ending the game immediately
  \end{enumerate}
  \item For cooperation to be sustainable, the cooperative payoff must exceed ...
  \begin{enumerate}
    \item the number of rounds
    \item the gain from a one-off defection followed by punishment
    \item the discount factor
    \item the total of all stage payoffs
  \end{enumerate}
  \item A larger discount factor \(\delta\) means ...
  \begin{enumerate}
    \item the future matters less
    \item the future matters more, making cooperation easier to sustain
    \item the game is shorter
    \item defection is always best
  \end{enumerate}
  \item The utility of an infinitely repeated game with discount factor \(\delta\) is ...
  \begin{enumerate}
    \item \(\max_i u(s_r(i), s_c(i))\)
    \item \(\sum_{i=0}^{\infty} \delta^i u(s_r(i), s_c(i))\)
    \item \(\sum_{i=0}^{\infty} u(s_r(i), s_c(i))\)
    \item \(\delta \, u(s_r(0), s_c(0))\)
  \end{enumerate}
  \item An individually rational payoff pair is one that ...
  \begin{enumerate}
    \item equals the temptation payoff
    \item gives one player everything
    \item minimises both players' payoffs
    \item exceeds the stage-game Nash equilibrium payoffs for both players
  \end{enumerate}
  \item For a patient enough population the Folk Theorem constructs a subgame perfect equilibrium achieving ...
  \begin{enumerate}
    \item any feasible individually rational payoff pair
    \item no payoff at all
    \item only the stage Nash payoff
    \item only the cooperative payoff
  \end{enumerate}
  \item In the Folk Theorem's trigger strategy, a deviation is punished by ...
  \begin{enumerate}
    \item leaving the game
    \item forgiving after one round
    \item raising the discount factor
    \item reverting to the stage-game Nash equilibrium for all future rounds
  \end{enumerate}
  \item Against a trigger strategy, the most tempting time to deviate is ...
  \begin{enumerate}
    \item a randomly chosen stage
    \item the first stage, since later gains are discounted more heavily
    \item never, deviation always pays
    \item the last stage
  \end{enumerate}
\end{enumerate}
\end{document}
